\section{实验部分}

\subsection{Ti${_3}$C${_2}$T${_x}$ MXene 的合成}

Ti$_{3}$C$_{2}$T$_{x}$ MXene 通过 LiF/HCl 刻蚀剂从 Ti$_{3}$AlC$_{2}$ MAX 相（纯度 $\geq$ 98\%，$\leq$ 400 目）中选择性刻蚀 Al 层制备。典型步骤为：称取 3.2 g LiF 溶于 40 mL HCl（12 mol L$^{-1}$）中，在聚四氟乙烯（PTFE）烧杯中室温磁力搅拌 30 min。随后，将 2 g Ti$_{3}$AlC$_{2}$ 粉末缓慢加入混合溶液中，在 40 $^{\circ}$C 恒温水浴中持续搅拌反应 48 h。将反应产物以 5,000 rpm 离心 1 min，弃去上清液，收集底部沉淀。用 1 mol L$^{-1}$ HCl 和超纯水对沉淀进行反复离心洗涤（5,000 rpm，1 min/次），直至上清液 pH 值接近中性。将所得 MXene 分散于 80 mL 超纯水中，在冰浴中超声处理 30 min。超声后，将分散液以 3,500 rpm 离心 15 min，收集上层深色液体，即为高浓度 MXene 胶体溶液（20--25 mg mL$^{-1}$）。最终所得 MXene 悬浮液通入 N$_{2}$ 保护，于冷藏条件下保存。

\subsection{Bi/MXene 复合材料的合成}

Bi/MXene 复合材料通过水热路线合成。首先，将冷冻干燥后的 MXene 气凝胶（0.21 g）加入含 2 mmol BiCl$_{3}\cdot x$H$_{2}$O 的 40 mL 无水乙醇中。混合物经磁力搅拌 3 h 以确保均匀分散，随后转移至聚四氟乙烯内衬的不锈钢水热反应釜中，在 180 $^{\circ}$C 下反应 16 h。此过程中，Bi$^{3+}$ 发生水解，在 MXene 表面成核形成 Bi(OH)$_{3}$ 纳米颗粒，得到 Bi(OH)$_{3}$/MXene 中间产物。反应结束后，产物通过离心收集，用无水乙醇和超纯水多次清洗，在真空烘箱中 60 $^{\circ}$C 干燥。最后，将干燥粉末置于管式炉中，在 Ar 气氛下 350 $^{\circ}$C 煅烧 2 h，使 Bi(OH)$_{3}$ 热分解并还原为金属 Bi，得到 Bi/MXene 复合材料。

作为对比，在不加入 MXene 的相同条件下制备了纯 Bi 样品。

\subsection{材料表征}

采用场发射扫描电子显微镜（SEM, JEOL JSM-6700F）观察样品的形貌和微观结构，并配备能量色散 X 射线光谱仪（EDS）进行元素分布分析。采用高分辨率透射电子显微镜（HRTEM, JEOL JEM-2010）对原子尺度界面结构进行表征。采用 X 射线衍射仪（XRD, Rigaku D/Max2500）以 Cu K$\alpha$ 辐射（$\lambda = 1.5406$ \AA）测定晶体物相。采用 X 射线光电子能谱仪（XPS, Thermo Fisher Scientific ESCALAB 250Xi）分析元素的化学状态。

\subsection{电化学测试}

电化学测试采用在充满 Ar 气的手套箱（H$_{2}$O, O$_{2}$ < 0.1 ppm）中组装的 CR2032 扣式半电池。工作电极的制备方法为：将活性材料（Bi/MXene 或纯 Bi）、Super P 导电炭黑、羧甲基纤维素钠（CMC）和丁苯橡胶（SBR）按 7:1:1:1 的质量比在超纯水中均匀混合成浆料，涂覆于铜箔集流体上。涂覆后的铜箔在真空烘箱中 60 $^{\circ}$C 干燥 12 h。活性材料的面载量约为 1 mg cm$^{-2}$。以金属钠箔同时作为对电极和参比电极，隔膜为玻璃纤维隔膜（Whatman GF/D），电解液为 1.0 mol L$^{-1}$ NaPF$_{6}$ 的乙二醇二甲醚（DME）溶液。

恒流充放电（GCD）测试在 Neware 电池测试系统上进行，电压区间为 0.01--1.8 V（vs.\ Na$^{+}$/Na）。倍率性能在 0.1 至 20 A g$^{-1}$ 的电流密度范围内评估。长循环稳定性在 2 A g$^{-1}$ 下进行 5,000 次循环测试。循环伏安（CV）测试在 Princeton Applied Research 电化学工作站上进行，扫速分别为 0.2、0.5、1.0、2.0 和 5.0 mV s$^{-1}$，电压窗口同上。电化学阻抗谱（EIS）测试的频率范围为 100 kHz 至 0.01 Hz，交流振幅为 5 mV。恒电流间歇滴定技术（GITT）测试采用 0.05 mA 的电流脉冲，脉冲时间为 10 min，弛豫时间为 60 min。
